\section{Limitations}
\label{sec:limitations}

The study is a purposive single-coder pilot. The corpus was selected to establish conceptual coverage, hard cases, and metric feasibility. Direct-support counts therefore measure documented recurrence inside the selected evidence register. They do not estimate forum prevalence, operational incidence, or industry frequency. The absence of a second coder means the study cannot report inter-rater agreement. Explicit code boundaries, difficult-case segmentation, counterexamples, executable value checks, and a claim ledger reduce ambiguity without replacing independent replication.

Public discourse creates selection and censoring. Shareable problems are overrepresented relative to confidential, regulated, or commercially sensitive incidents. Advanced platforms and active communities may generate more discussion because they have larger installed bases, more complex uses, or stronger public participation. A thread can end after a direct message, vendor support exchange, meeting, Slack migration, or local fix. Public open status therefore cannot be used as a general non-resolution indicator.

The field-level metrics have small and heterogeneous denominators. Their purpose is to show that bounded technical measurement is possible. Means across device interfaces, deployment objects, resource definitions, and test claims are not a common performance scale. Some component fields may be inapplicable to particular systems. Unknown is excluded from known-component means, which prevents false absence and can raise uncertainty when public evidence is sparse.

Association measures arise from multi-label codes in a purposive corpus. Threads are not independent draws, and technical categories influence which constructs are discussed together. Phi coefficients and lift identify relationships for qualitative adjudication and prospective testing. They do not estimate causal effects. The B2/B10 matched-case relationship has additional measurement overlap and is used only as convergent validity.

Source-reported product results remain source-reported. The VerisFlow rate is preserved because its numerator, denominator, random-worklist description, simulation environment, and trace-matching criterion are explicit in the forum post. The paper does not independently replicate that experiment. Its contribution is to bind the number to simulation-trace agreement and exclude deeper claims about dry, wet, assay, or production performance.

The taxonomy should remain revisable. The ten constructs reflect a specialist public community over a four-year period. Future work may identify new constructs, split current constructs, or show that some boundaries merge under standardized infrastructures. The public repository exposes schemas, data, tests, and a correction pathway so that revisions can be tracked without silently changing prior releases.
